% Part: proof-theory
% Chapter: normalization
% Section: translations

\documentclass[../../../include/open-logic-section]{subfiles}

\begin{document}

\olfileid{pt}{nor}{tr}

\olsection{في النقل بين \usetoken{P}{proof} السوية في~\Log{N2i} و\Log{G2i}}

للـ!!^{proof}s في~\Log{N2i} نقل إلى~!!{proof}s في~\Log{G2i}، وللثانية
نقل إلى الأولى. وأول ما يظهر أن نقل~!!{proof}s الخالية من القطع في
\Log{G2i} لا ينتج إلا~!!{proof}s سوية في~\Log{N2i}.

وقبل صياغة النتيجة نثبت اصطلاحًا يصل مقدمات المتتاليات في~\Log{N2i}
بمقدماتها في~\Log{G2i}. فنقول إن مجموعة~${\OLId{greek-variable}{Gamma}}'$ من~!!{formula}s
الموسومة \emph{تقابل} المجموعة المتعددة~${\OLId{greek-variable}{Gamma}}$ من~!!{formula}s غير
الموسومة، إذا كان، لكل~!!{formula}~$!A$، عدد~!!{element}s~${\OLId{greek-variable}{Gamma}}'$ التي
على الصورة~${\OLId{latin-ordinary}{x}}:!A$ غير زائد على تعددية~$!A$ في~${\OLId{greek-variable}{Gamma}}$.

\begin{cor}
إذا كان $\Log{G2i} \Proves {\OLId{greek-variable}{Gamma}} \Sequent {\OLId{greek-variable}{Delta}}$، فثمة
!!{proof} سوي~$\delta$ في~\Log{N2i} للمتتالية~${\OLId{greek-variable}{Gamma}}' \Sequent !C$، حيث
$!C\ident\lfalse$ إذا كانت~${\OLId{greek-variable}{Delta}}$ خالية، و$!C\ident !A$ إذا كانت
${\OLId{greek-variable}{Delta}}=\{!A\}$، وحيث تقابل مجموعة~${\OLId{greek-variable}{Gamma}}'$ من~!!{formula}s الموسومة
المجموعة المتعددة~${\OLId{greek-variable}{Gamma}}$؛ أي إنه، لكل~!!{formula}~$!A$، يكون عدد
!!{element}s المجموعة~${\OLId{greek-variable}{Gamma}}'$ التي صورتها~${\OLId{latin-ordinary}{x}}:!A$ أصغر من أو يساوي
تعددية~$!A$ (أي عدد نسخ~$!A$) في~${\OLId{greek-variable}{Gamma}}$.
\end{cor}

\begin{proof}
وبيانه في فحص برهان~\olref[nat][tgi]{prop:G2i-to-N2i}: فإننا نلحق قواعد
!!{introduction} بنهايات~!!{proof}s، ولا نضع قواعد~!!{elimination} إلا
في أعالي~!!{proof}s؛ فلا ندخل قط قاعدة~!!{introduction} فوق
!!a{elimination}.
\end{proof}

إلا أن نقل قاعدة~\Cut{} المبين في~\olref[nat][tgi]{prop:G2i-to-N2i}
قد يحدث قطوعًا في~!!{proof}~$\delta$ من~\Log{N2i}. فإذا كان آخر~!!{proof}
في~\Log{G2i} بـ~\Cut{} وكان~!!{proof}sه الجزئيان المباشران~$\pi_{\OLMathNumeral{1}}$ و$\pi_{\OLMathNumeral{2}}$
للمتتاليتين~${\OLId{greek-variable}{Gamma}}_{\OLMathNumeral{1}} \Sequent !A$ و
$!A,{\OLId{greek-variable}{Gamma}}_{\OLMathNumeral{2}} \Sequent {\OLId{greek-variable}{Delta}}_{\OLMathNumeral{2}}$ على الترتيب، فإننا نوصل ترجمة~$\pi_{\OLMathNumeral{1}}$،
وهي~$\delta_{\OLMathNumeral{1}}$، بمسلّمات~${\OLId{latin-ordinary}{x}}:!A \Sequent !A$ الواقعة في ترجمة~$\pi_{\OLMathNumeral{2}}$،
وهي~$\delta_{\OLMathNumeral{2}}$. فإذا انتهى~$\delta_{\OLMathNumeral{1}}$ بقاعدة~!!a{introduction} صيغتها
الرئيسة~!!{formula}~$!A$، وكانت المسلّمة~${\OLId{latin-ordinary}{x}}:!A \Sequent !A$ في~$\delta_{\OLMathNumeral{2}}$
المقدمة الكبرى لقاعدة~!!a{elimination}، نتج عن ذلك قطع في~$\delta$.

وفي الجهة الأخرى تنقل~!!{proof}s في~\Log{N2i} إلى~!!{proof}s في
\Log{G2i}. ففي برهان~\olref[nat][tni]{prop:N2-to-G2} ينشأ عن نقل~\Elim{\land} و
\Elim{\lif} و\Elim{\lnot} و\Elim{\lforall} استدلالات~\Cut{} في
!!{proof}~\Log{G2i} الناتج. لكن ذلك يجتنب متى كان~!!{proof}~\Log{N2i}
الذي نبدأ به سويًا.

ونسمي قائمة~${\OLId{latin-looped}{S}}_{\OLMathNumeral{1}}$، \dots، ${\OLId{latin-looped}{S}}_{\OLId{latin-ordinary}{n}}$ من وقائع المتتاليات في~!!{proof}
$\delta$ من~\Log{N2i} \emph{فرعًا} إذا كانت~${\OLId{latin-looped}{S}}_{\OLMathNumeral{1}}$ مسلّمة، وكان~${\OLId{latin-looped}{S}}_{{\OLId{latin-ordinary}{i}}+{\OLMathNumeral{1}}}$
نتيجة الاستدلال متى كان~${\OLId{latin-looped}{S}}_{\OLId{latin-ordinary}{i}}$ مقدمته الكبرى، ولم تكن العقدة الأخيرة
${\OLId{latin-looped}{S}}_{\OLId{latin-ordinary}{n}}$ مقدمة كبرى لاستدلال. فالفرع يتعقب المتتاليات هبوطًا في~$\delta$
ابتداء من المسلّمات، ويقف عند المقدمات الصغرى لقواعد~!!{elimination}.
ونسمي \emph{الفرع الرئيس} لـ~$\delta$ كل فرع تكون~${\OLId{latin-looped}{S}}_{\OLId{latin-ordinary}{n}}$ فيه المتتالية
النهائية. ولا يخلو~!!{proof}~$\delta$ من فرع رئيس واحد على الأقل؛ لأن
لكل قاعدة مقدمة كبرى، ولا تكون لها مقدمتان كبريان إلا~\Intro{\land}.

\begin{prop}\ollabel{prop:normal-N2i-to-G2i}
إذا كان~$\delta$ !!{proof} سويًا في~$\Log{N2i}$ للمتتالية
${\OLId{greek-variable}{Gamma}} \Sequent !A$، فثمة~!!{proof}~$\pi$ خالٍ من القطع، وهو
!!{proof} في~$\Log{G2i}$
للمتتالية~${\OLId{greek-variable}{Gamma}}' \Sequent {\OLId{greek-variable}{Delta}}$، حيث~${\OLId{greek-variable}{Gamma}}'$ هي المجموعة المتعددة
من~!!{formula}s الناتجة من~${\OLId{greek-variable}{Gamma}}$ بحذف الوسوم، و
${\OLId{greek-variable}{Delta}}=\{!A\}$ إذا لم تكن~$!A$ هي~$\lfalse$، و
${\OLId{greek-variable}{Delta}}=\emptyset$ إذا كانت إياها.
\end{prop}

\begin{proof}
ونسلك هذه المرة الاستقراء على عدد الاستدلالات في~!!{proof}~$\delta$
للمتتالية~${\OLId{greek-variable}{Gamma}} \Sequent !A$. فإن خلا~$\delta$ من الاستدلال كان
المسلّمة~${\OLId{latin-ordinary}{x}}:!A \Sequent !A$، وأخذنا~$\pi$ المسلّمة المقابلة
$!A \Sequent !A$، أو~$\lfalse \Sequent \ $ إذا كان
$!A\ident\lfalse$.

وإذا كان آخر~$\delta$ قاعدة استدلال، قسمنا إلى الحالات الآتية:
\begin{enumerate}
  \item إذا كانت~${\OLId{latin-looped}{R}}$ قاعدة~!!a{introduction} أو~\FalseInt، جرت الترجمة
  كما في برهان~\olref[nat][tni]{prop:N2-to-G2} تمامًا، ولم تحتج إلى~\Cut.

  \item إذا كانت آخر قاعدة~${\OLId{latin-looped}{R}}$ في~$\delta$ قاعدة~!!a{elimination}، فلاحظ
  ما يأتي:
  \begin{enumerate}
    \item لا يمر فرع رئيس في~$\delta$ بأي قواعد~!!{introduction}، وإلا
    لكانت قاعدة~!!a{introduction} أو~\FalseInt{} متبوعةً بقاعدة
    !!a{elimination} على الفرع الرئيس، ولما كان~$\delta$ سويًا.
    \item إذ إن الفروع الرئيسة في~$\delta$ لا تمر بقواعد~!!{introduction}،
    فلا يوجد إلا فرع رئيس واحد. (لقاعدة~\Intro{\land} وحدها مقدمتان
    كبريان، ولذلك يجب أن تمر الفروع الرئيسة المتعددة بمقدمتي قاعدة من هذا
    النوع.)
    \item إذا اشتمل الفرع الرئيس على المقدمة الكبرى لـ~\Elim{\lor} أو
    \Elim{\lexists}، فلا توجد إلا قاعدة واحدة من هذا النوع وتكون الأخيرة،
    أي~${\OLId{latin-looped}{R}}$. (وإلا لاحتوى الفرع الرئيس قاعدة~\Elim{\lor} أو
    \Elim{\lexists} تكون نتيجتها مقدمةً كبرى لقاعدة~!!a{elimination}، ولما
    كان البرهان سويًا، إذ يمكن عندئذ تطبيق تحويل تبديل.)
    \item لا تفرّغ أي~!!{elimination} غير~\Elim{\lor} و\Elim{\lexists}
    فروضًا، أي لا تغيّر السياق الأيسر للمتتاليات. ولا تفرّغ هاتان القاعدتان
    إلا فروضًا فوق مقدمة صغرى، أي لا تغيّران السياق الأيسر إلا لمقدمتيهما
    الصغريين. وبعبارة أخرى، إذا كان سياق متتالية~${\OLId{latin-looped}{S}}_{\OLId{latin-ordinary}{i}}$ على الفرع الرئيس
    لـ~$\delta$ هو~${\OLId{greek-variable}{Gamma}}_{\OLMathNumeral{1}}$، فإن سياق النتيجة~${\OLId{latin-looped}{S}}_{{\OLId{latin-ordinary}{i}}+{\OLMathNumeral{1}}}$ للاستدلال الذي
    تكون~${\OLId{latin-looped}{S}}_{\OLId{latin-ordinary}{i}}$ مقدمته الكبرى هو سياق حديث~${\OLId{greek-variable}{Gamma}}_{\OLMathNumeral{2}} \supseteq {\OLId{greek-variable}{Gamma}}_{\OLMathNumeral{1}}$.
    \item ونتيجةً لذلك، إذا كانت~${\OLId{latin-ordinary}{x}}:!C \Sequent !C$ المتتالية العليا
    ${\OLId{latin-looped}{S}}_{\OLMathNumeral{1}}$ في الفرع الرئيس، فإن~${\OLId{latin-ordinary}{x}}:!C \in {\OLId{greek-variable}{Gamma}}$.
  \end{enumerate}
  والآن نميّز الحالات بحسب صورة~$!C$. وبما أن كل~${\OLId{latin-looped}{S}}_{\OLId{latin-ordinary}{i}}$ في الفرع الرئيس
  مقدمة كبرى لقاعدة~!!a{elimination}، ينطبق ذلك كذلك على
  ${\OLId{latin-ordinary}{x}}:!C \Sequent !C$.

  \item إذا كان $!C \ident !D \land !E$، كان~$\delta$ بالصورة:
  \[
  \Axiom${\OLId{latin-ordinary}{x}}:!D \land !E \fCenter !D \land !E$
  \RightLabel{\Elim{\land}[{\OLMathNumeral{1}}]}
  \UnaryInf${\OLId{latin-ordinary}{x}}:!D \land !E \fCenter !D$
  \Deduce${\OLId{latin-ordinary}{x}}:!D \land !E, {\OLId{greek-variable}{Gamma}}_{\OLMathNumeral{1}} \fCenter !A$
  \DisplayProof
  \]
  لا يمر الفرع الرئيس، الذي يشتمل على~${\OLId{latin-ordinary}{x}}:!D \land !E \Sequent !D$، بنتيجة
  قاعدة~\Intro{\lif} أو~\Intro{\lnot}، ولا بمقدمة صغرى لـ~\Elim{\lor} أو
  \Elim{\lexists}؛ ومن ثم تبقى~!!{formula}s السياق في
  ${\OLId{latin-ordinary}{x}}:!D \land !E$ بلا تغيير على طول الفرع الرئيس. وهذا يفسر وجوب احتواء
  سياق المتتالية النهائية على~${\OLId{latin-ordinary}{x}}:!D \land !E$. وفوق ذلك يمكننا تحويل
  !!{proof}~$\delta$ إلى~!!{proof}~$\delta_{\OLMathNumeral{1}}$ باستبدال الـ~!!{proof}
  الجزئي المنتهي بـ~\Elim{\land} بالمسلمة~${\OLId{latin-ordinary}{x}}:!D \Sequent !D$، واستبدال
  ${\OLId{latin-ordinary}{x}}:!D \land !E$ بـ~${\OLId{latin-ordinary}{x}}:!D$ في سياق كل متتالية على الفرع الرئيس؛ أي نحوّل
  \[
  \bottomAlignProof
  \Axiom${\OLId{latin-ordinary}{x}}:!D \land !E \fCenter !D \land !E$
  \RightLabel{\Elim{\land}[{\OLMathNumeral{1}}]}
  \UnaryInf${\OLId{latin-ordinary}{x}}:!D \land !E \fCenter !D$
  \Deduce${\OLId{latin-ordinary}{x}}:!D \land !E, {\OLId{greek-variable}{Gamma}}_{\OLMathNumeral{1}} \fCenter !A$
  \DisplayProof
  \quad\text{إلى}\quad
  \bottomAlignProof
  \Axiom${\OLId{latin-ordinary}{x}}:!D \fCenter !D$
  \Deduce${\OLId{latin-ordinary}{x}}:!D, {\OLId{greek-variable}{Gamma}}_{\OLMathNumeral{1}} \fCenter !A$
  \DisplayProof
  \]
  ينطبق فرض الاستقراء على~$\delta_{\OLMathNumeral{1}}$، لأن عدد الاستدلالات في~$\delta$
  يساوي عدد الاستدلالات في~$\delta_{\OLMathNumeral{1}}$ زائد~${\OLMathNumeral{1}}$.

  وبفرض الاستقراء يوجد~!!{proof}~$\pi_{\OLMathNumeral{1}}$ في~\Log{G2i} للمتتالية
  $!D,{\OLId{greek-variable}{Gamma}}_{\OLMathNumeral{1}}' \Sequent {\OLId{greek-variable}{Delta}}$. ونحصل على~!!{proof}~$\pi$ بتطبيق
  \LeftR{\land}، مع إبقاء الطرف الأيمن~${\OLId{greek-variable}{Delta}}$ نفسه:
  \[
  \AxiomC{}
  \RightLabel{$\pi_{\OLMathNumeral{1}}$}
  \Deduce$!D, {\OLId{greek-variable}{Gamma}}_{\OLMathNumeral{1}}' \fCenter {\OLId{greek-variable}{Delta}}$
  \RightLabel{\LeftR{\land}}
  \UnaryInf$!D \land !E, {\OLId{greek-variable}{Gamma}}_{\OLMathNumeral{1}}' \fCenter {\OLId{greek-variable}{Delta}}$
  \DisplayProof
  \]
  وإذا كانت~${\OLId{greek-variable}{Delta}}$ خالية فلا نطبّق إضعافًا يمينيًا؛ إذ يجب أن تبقى خالية.

  \item إذا كان $!C \ident !D \lor !E$، فلا بد أن يكون المقدمة الكبرى لـ
  \Elim{\lor}، ولا بد أن يكون هذا الاستدلال آخر استدلال~${\OLId{latin-looped}{R}}$ على الفرع
  الرئيس. أي إن~$\delta$ بالصورة:
  \[
  \Axiom${\OLId{latin-ordinary}{x}}:!D \lor !E \fCenter !D \lor !E$
  \AxiomC{}
  \RightLabel{$\delta_{\OLMathNumeral{1}}$}
  \Deduce${\OLId{latin-ordinary}{y}}:!D, {\OLId{greek-variable}{Gamma}}_{\OLMathNumeral{1}} \fCenter!A$
  \AxiomC{}
  \RightLabel{$\delta_{\OLMathNumeral{2}}$}
  \Deduce${\OLId{latin-ordinary}{z}}:!E, {\OLId{greek-variable}{Gamma}}_{\OLMathNumeral{2}} \fCenter !A$
  \RightLabel{\Elim{\lor}}
  \TrinaryInf${\OLId{latin-ordinary}{x}}:!D \lor !E, {\OLId{greek-variable}{Gamma}}_{\OLMathNumeral{1}}, {\OLId{greek-variable}{Gamma}}_{\OLMathNumeral{2}} \fCenter !A$
  \DisplayProof
  \]
  ينطبق فرض الاستقراء على~!!{proof}s~$\delta_{\OLMathNumeral{1}}$ و$\delta_{\OLMathNumeral{2}}$؛ فنحصل على
  !!{proof}s~$\pi_{\OLMathNumeral{1}}$ و$\pi_{\OLMathNumeral{2}}$ في~\Log{G2i} للمتتاليتين
  $!D,{\OLId{greek-variable}{Gamma}}_{\OLMathNumeral{1}}' \Sequent {\OLId{greek-variable}{Delta}}$ و$!E,{\OLId{greek-variable}{Gamma}}_{\OLMathNumeral{2}}' \Sequent {\OLId{greek-variable}{Delta}}$.
  ونطبّق~$\LeftR{\lor}$ للحصول على~$\pi$:
  \[
  \AxiomC{}
  \RightLabel{$\pi_{\OLMathNumeral{1}}$}
  \Deduce$!D, {\OLId{greek-variable}{Gamma}}_{\OLMathNumeral{1}}' \fCenter {\OLId{greek-variable}{Delta}}$
  \AxiomC{}
  \RightLabel{$\pi_{\OLMathNumeral{2}}$}
  \Deduce$!E, {\OLId{greek-variable}{Gamma}}_{\OLMathNumeral{2}}' \fCenter {\OLId{greek-variable}{Delta}}$
  \RightLabel{\LeftR{\lor}}
  \BinaryInf$!D \lor !E, {\OLId{greek-variable}{Gamma}}_{\OLMathNumeral{1}}', {\OLId{greek-variable}{Gamma}}_{\OLMathNumeral{2}}' \fCenter {\OLId{greek-variable}{Delta}}$
  \DisplayProof
  \]
  وإذا لم تفرّغ~$\Elim{\lor}$ الفرض~${\OLId{latin-ordinary}{y}}:!D$ أو~${\OLId{latin-ordinary}{z}}:!E$ (أي إذا لم يكن
  موجودًا في سياق المقدمة الصغرى)، أضفنا ما يلزم فقط من تطبيقات
  \LeftR{\Weakening}. ولا نضيف إضعافًا يمينيًا عندما تكون~${\OLId{greek-variable}{Delta}}$ خالية؛
  فمثلًا، إذا غاب الفرضان:
  \[
  \AxiomC{}
  \RightLabel{$\pi_{\OLMathNumeral{1}}$}
  \Deduce${\OLId{greek-variable}{Gamma}}_{\OLMathNumeral{1}}' \fCenter {\OLId{greek-variable}{Delta}}$
  \RightLabel{\LeftR{\Weakening}}
  \UnaryInf$!D, {\OLId{greek-variable}{Gamma}}_{\OLMathNumeral{1}}' \fCenter {\OLId{greek-variable}{Delta}}$
  \AxiomC{}
  \RightLabel{$\pi_{\OLMathNumeral{2}}$}
  \Deduce${\OLId{greek-variable}{Gamma}}_{\OLMathNumeral{2}}' \fCenter {\OLId{greek-variable}{Delta}}$
  \RightLabel{\LeftR{\Weakening}}
  \UnaryInf$!E, {\OLId{greek-variable}{Gamma}}_{\OLMathNumeral{2}}' \fCenter {\OLId{greek-variable}{Delta}}$
  \RightLabel{\LeftR{\lor}}
  \BinaryInf$!D \lor !E, {\OLId{greek-variable}{Gamma}}_{\OLMathNumeral{1}}', {\OLId{greek-variable}{Gamma}}_{\OLMathNumeral{2}}' \fCenter {\OLId{greek-variable}{Delta}}$
  \DisplayProof
  \]

  وفي كلتا الصورتين قد يحتوي المقدم المعروض نسخًا زائدة إذا وقع الفرض
  الموسوم نفسه في أكثر من واحد من السياقات
  $\{{\OLId{latin-ordinary}{x}}:!D \lor !E\}$ و${\OLId{greek-variable}{Gamma}}_{\OLMathNumeral{1}}$ و${\OLId{greek-variable}{Gamma}}_{\OLMathNumeral{2}}$. فنلحق ما يلزم من
  تطبيقات~\LeftR{\Contraction}، ولا نحذف نسخة ناشئة من فرض موسوم آخر وإن
  اتفقت الصيغتان، حتى تصير المتتالية النهائية بالضبط
  \[
  {\OLId{greek-variable}{Gamma}}' \Sequent {\OLId{greek-variable}{Delta}},
  \]
  حيث~${\OLId{greek-variable}{Gamma}}'$ هي المجموعة المتعددة الناتجة من سياق نتيجة~\Elim{\lor}
  بحذف الوسوم.

  \item إذا كان $!C \ident !D \lif !E$، كان~$\delta$ بالصورة:
  \[
  \Axiom${\OLId{latin-ordinary}{x}}:!D \lif !E \fCenter !D \lif !E$
  \AxiomC{}
  \RightLabel{$\delta_{\OLMathNumeral{1}}$}
  \Deduce${\OLId{greek-variable}{Gamma}}_{\OLMathNumeral{1}} \fCenter!D$
  \RightLabel{\Elim{\lif}}
  \BinaryInf${\OLId{latin-ordinary}{x}}:!D \lif !E, {\OLId{greek-variable}{Gamma}}_{\OLMathNumeral{1}} \fCenter !E$
  \RightLabel{$\delta_{\OLMathNumeral{2}}$}
  \Deduce${\OLId{latin-ordinary}{x}}:!D \lif !E, {\OLId{greek-variable}{Gamma}}_{\OLMathNumeral{1}}, {\OLId{greek-variable}{Gamma}}_{\OLMathNumeral{2}} \fCenter !A$
  \DisplayProof
  \]
  لاحظ هنا أن الفرع الرئيس، الذي يشتمل على
  ${\OLId{latin-ordinary}{x}}:!D \lif !E,{\OLId{greek-variable}{Gamma}}_{\OLMathNumeral{1}} \Sequent !E$، لا يمر بنتيجة~\Intro{\lif} أو
  \Intro{\lnot} ولا بمقدمة صغرى لـ~\Elim{\lor} أو~\Elim{\lexists}؛ ولذلك
  تبقى~!!{formula}s السياق في~${\OLId{latin-ordinary}{x}}:!D \lif !E,{\OLId{greek-variable}{Gamma}}_{\OLMathNumeral{1}}$ بلا تغيير على طول
  الفرع الرئيس. وهذا يفسر وجوب احتواء سياق المتتالية النهائية على
  ${\OLId{latin-ordinary}{x}}:!D \lif !E,{\OLId{greek-variable}{Gamma}}_{\OLMathNumeral{1}}$. وفوق ذلك يمكننا تحويل قطعة~!!{proof}~$\delta_{\OLMathNumeral{2}}$
  إلى~!!{proof} صحيح~$\delta_{\OLMathNumeral{2}}'$ باستبدال الـ~!!{proof} الجزئي المنتهي
  بـ~\Elim{\lif} بالمسلمة~${\OLId{latin-ordinary}{x}}:!E \Sequent !E$، وباستبدال
  ${\OLId{latin-ordinary}{x}}:!D \lif !E,{\OLId{greek-variable}{Gamma}}_{\OLMathNumeral{1}}$ في سياق كل متتالية على الفرع الرئيس بـ~${\OLId{latin-ordinary}{x}}:!E$؛
  أي نحوّل
  \[
  \bottomAlignProof
  \Axiom${\OLId{latin-ordinary}{x}}:!D \lif !E, {\OLId{greek-variable}{Gamma}}_{\OLMathNumeral{1}} \fCenter !E$
  \RightLabel{$\delta_{\OLMathNumeral{2}}$}
  \Deduce${\OLId{latin-ordinary}{x}}:!D \lif !E, {\OLId{greek-variable}{Gamma}}_{\OLMathNumeral{1}}, {\OLId{greek-variable}{Gamma}}_{\OLMathNumeral{2}} \fCenter !A$
  \DisplayProof
  \quad\text{إلى}\quad
  \bottomAlignProof
  \Axiom${\OLId{latin-ordinary}{x}}:!E \fCenter !E$
  \RightLabel{$\delta_{\OLMathNumeral{2}}'$}
  \Deduce${\OLId{latin-ordinary}{x}}:!E, {\OLId{greek-variable}{Gamma}}_{\OLMathNumeral{2}} \fCenter !A$
  \DisplayProof
  \]
  ينطبق فرض الاستقراء على~$\delta_{\OLMathNumeral{1}}$ و$\delta_{\OLMathNumeral{2}}'$، لأن عدد الاستدلالات في
  $\delta$ هو مجموع عددي الاستدلالات في~$\delta_{\OLMathNumeral{1}}$ و$\delta_{\OLMathNumeral{2}}$ زائد~${\OLMathNumeral{1}}$
  لاستدلال~\Elim{\lif} في بداية الفرع الرئيس. وإذا كان هذا الاستدلال نفسه
  آخر استدلال~${\OLId{latin-looped}{R}}$ في~$\delta$، احتوى~$\delta_{\OLMathNumeral{2}}$ صفرًا من الاستدلالات.

  وبفرض الاستقراء توجد~!!{proof}s~$\pi_{\OLMathNumeral{1}}$ و$\pi_{\OLMathNumeral{2}}$ في~\Log{G2i}، الأولى
  للمتتالية~${\OLId{greek-variable}{Gamma}}_{\OLMathNumeral{1}}' \Sequent !D$ إذا كان~$!D\not\ident\lfalse$، أو
  للطرف الأيمن الخالي إذا كان~$!D\ident\lfalse$، والثانية للمتتالية
  $!E,{\OLId{greek-variable}{Gamma}}_{\OLMathNumeral{2}}' \Sequent {\OLId{greek-variable}{Delta}}$. ونحصل على~!!{proof}~$\pi$ بتطبيق
  \LeftR{\lif}:
  \[
  \AxiomC{}
  \RightLabel{$\pi_{\OLMathNumeral{1}}$}
  \Deduce${\OLId{greek-variable}{Gamma}}_{\OLMathNumeral{1}}' \fCenter !D$
  \AxiomC{}
  \RightLabel{$\pi_{\OLMathNumeral{2}}$}
  \Deduce$!E, {\OLId{greek-variable}{Gamma}}_{\OLMathNumeral{2}}' \fCenter {\OLId{greek-variable}{Delta}}$
  \RightLabel{\LeftR{\lif}}
  \BinaryInf$!D \lif !E, {\OLId{greek-variable}{Gamma}}_{\OLMathNumeral{1}}', {\OLId{greek-variable}{Gamma}}_{\OLMathNumeral{2}}' \fCenter {\OLId{greek-variable}{Delta}}$
  \DisplayProof
  \]
  وإذا كان~$!D\ident\lfalse$ وحده، نطبّق~\RightR{\Weakening} على~$\pi_{\OLMathNumeral{1}}$
  لتوفير المقدمة الأولى لـ~\LeftR{\lif}. ولا نطبّق إضعافًا يمينيًا على
  $\pi_{\OLMathNumeral{2}}$ أو على النتيجة إذا كانت~$!A\ident\lfalse$؛ فتبقى~${\OLId{greek-variable}{Delta}}$ خالية:
  \[
  \AxiomC{}
  \RightLabel{$\pi_{\OLMathNumeral{1}}$}
  \Deduce${\OLId{greek-variable}{Gamma}}_{\OLMathNumeral{1}}' \fCenter $
  \RightLabel{\RightR{\Weakening}}
  \UnaryInf${\OLId{greek-variable}{Gamma}}_{\OLMathNumeral{1}}' \fCenter !D$
  \AxiomC{}
  \RightLabel{$\pi_{\OLMathNumeral{2}}$}
  \Deduce$!E, {\OLId{greek-variable}{Gamma}}_{\OLMathNumeral{2}}' \fCenter {\OLId{greek-variable}{Delta}}$
  \RightLabel{\LeftR{\lif}}
  \BinaryInf$!D \lif !E, {\OLId{greek-variable}{Gamma}}_{\OLMathNumeral{1}}', {\OLId{greek-variable}{Gamma}}_{\OLMathNumeral{2}}' \fCenter {\OLId{greek-variable}{Delta}}$
  \DisplayProof
  \]

  وهنا أيضًا قد يحتوي المقدم المعروض نسخًا زائدة إذا وقع الفرض الموسوم
  نفسه في أكثر من واحد من السياقات
  $\{{\OLId{latin-ordinary}{x}}:!D \lif !E\}$ و${\OLId{greek-variable}{Gamma}}_{\OLMathNumeral{1}}$ و${\OLId{greek-variable}{Gamma}}_{\OLMathNumeral{2}}$. فنلحق ما يلزم من
  تطبيقات~\LeftR{\Contraction}، مع إبقاء النسخ الناشئة من فروض مختلفة
  الوسوم وإن اتفقت صيغها، حتى تصير المتتالية النهائية بالضبط
  \[
  {\OLId{greek-variable}{Gamma}}' \Sequent {\OLId{greek-variable}{Delta}}.
  \]

\end{enumerate}
وعلى هذا القياس تعالج الحالات الباقية، وقد تركت تمارين. ولا تتناول الحجة
الحالة الكلاسيكية~\FalseCl؛ فمن ثم تختص المبرهنة بـ\Log{N2i} و\Log{G2i}.
\end{proof}

\begin{cor}
كل~!!{formula} يقع في~!!{proof} سوي من~\Log{N2i} هو~!!{formula} فرعية
لإحدى~!!{formula}s في المتتالية النهائية. وكل صيغة تقع في برهان سوي من
\Log{N1i} هي صيغة فرعية للصيغة النهائية أو لأحد الفروض المفتوحة.
\end{cor}

\begin{proof}
  من~\olref{prop:normal-N2i-to-G2i} ينقل كل~!!{proof} سوي في~\Log{N2i}
  إلى~!!{proof} خال من القطع في~\Log{G2i}. وبفحص البرهان نجد أن كل
  !!{formula} واقعة في~!!{proof}~\Log{N2i} واقعة أيضًا في~!!{proof}
  \Log{G2i}. ولما كان~\Log{G2i} خاليًا من قاعدة~\Cut{}، ثبتت له خاصية
  !!{formula} الفرعية. وأما نقل~\Log{N1i} إلى~\Log{N2i} فيحفظ السوية،
  ويثبت الفروض المفتوحة في مقدم المتتالية النهائية؛ فتنتج الحالة الثانية
  من الأولى.
\end{proof}

\end{document}
